\documentclass[12pt]{article}
\usepackage[utf8]{inputenc}
\usepackage[english]{babel}
\usepackage[T1]{fontenc} % accepting different input for PDF

% \usepackage{times}
\usepackage{newtxtext,newtxmath}
% \usepackage[document]{ragged2e}
% \usepackage{amssymb}
\usepackage{amsmath}
\usepackage{graphicx}
\usepackage{setspace}
\usepackage{csquotes}
\doublespacing
% \let\oldtabular\tabular
% \renewcommand{\tabular}{\setstretch{1}\oldtabular}
\usepackage[a4paper,width=150mm,top=25mm,bottom=25mm,bindingoffset=10mm]{geometry}
\usepackage[backend=biber, style=apa]{biblatex}
\usepackage{titlesec}
\usepackage{hyperref}
\usepackage{indentfirst}
\usepackage[skip = -25pt, indent=.5in]{parskip}
\graphicspath{CiHB/img/} % image file
\usepackage{pdfpages} % enabling input of pdf files by selecting specific pages
\usepackage{soul,color}
% \usepackage{longtable}
% \usepackage{colortbl,multirow,hhline}
\usepackage{tabularx}
\usepackage{ltablex}
\usepackage{pdflscape}
\usepackage{colortbl}
\usepackage{rotating}
\usepackage{tablefootnote}
\usepackage{tabularray}
% \DefTblrTemplate{firsthead,middlehead,lasthead}{default}{
% }
\DefTblrTemplate{firsthead}{default}{}
\DefTblrTemplate{firstfoot}{default}{
  \UseTblrTemplate{contfoot}{default}
  \UseTblrTemplate{caption}{default}
}
\DefTblrTemplate{lastfoot}{default}{
  \UseTblrTemplate{caption}{default}
}
% \DefTblrTemplate{middlefoot}{default}{
%   \UseTblrTemplate{contfoot}{default}
%   \UseTblrTemplate{capcont}{default}
% }
% \DefTblrTemplate{lastfoot}{default}{
%   \UseTblrTemplate{note}{default}
%   \UseTblrTemplate{remark}{default}
%   \UseTblrTemplate{capcont}{default}
% } 
\UseTblrLibrary{booktabs, color, longtblr, cell-span}
\usepackage[normalem]{ulem} % for multiple pages tables
\useunder{\uline}{\ul}{} % for multiple pages tables
\usepackage{siunitx}
% \UseTblrLibrary{varwidth}
\usepackage{changepage}
\usepackage[normalem]{ulem} % for multiple pages tables
\useunder{\uline}{\ul}{} % for multiple pages tables
\usepackage{float} % enabling tables' placement to be manually determined
% \usepackage{ragged2e}
\usepackage{makecell}
\setcellgapes{0pt}
\usepackage{comment}
\usepackage{tikz}
\NeedsTeXFormat{LaTeX2e}
\usepackage{venndiagram}[2018/06/07 v1.2 (NLCT) Venn diagrams]
\usepackage{subcaption}
\usetikzlibrary{backgrounds,positioning,calc,arrows.meta}
\usepackage{booktabs,array}

\RequirePackage{xkeyval}
\RequirePackage{etoolbox}


\renewcommand\cellalign{tl}   % top-left alignment inside \makecell
\renewcommand\theadalign{tl}  % top-left alignment for \thead if used

\usepackage{tikz}
\newcommand{\Closed}{\tikz[baseline=-0.6ex]\fill(0,0) circle (0.6ex);} % filled
\newcommand{\OpenClosed}{\tikz[baseline=-0.6ex]\fill[gray!50, draw=black](0,0) circle (0.6ex);} % gray
\newcommand{\Open}{\tikz[baseline=-0.6ex]\draw(0,0) circle (0.6ex);} % empty


% Title setup
\titleformat{\section}{\normalfont\Large\bfseries}{\thesection}{1em}{}
\titleformat{\subsection}{\normalfont\large\bfseries}{\thesubsection}{1em}{}
\titleformat{\subsubsection}{\normalfont\normalsize\bfseries}{\thesubsubsection}{1em}{}

% Bibliography style
\addbibresource{references.bib}

\begin{document}

% Title
\title{Envisioning Quantitative Resilience in the Age of the Great Acceleration}

\author{Sébastien Gillard \\
  Geneva School of Economics and Management, University of Geneva \\
  Email: sebastien.gillard@etu.unige.ch
  \and
  Lorena Rota \\
  xxxx \\
  xxxxx
  \and
  Thomas Maillart \\
  Geneva School of Economics and Management, University of Geneva \\
  Email: thomas.maillart@unige.ch \vspace{20pt}}

% alain et lorena

\date{\today}

\maketitle

% Abstract
\begin{abstract}
\vspace*{4mm}
Resilience has become a central concern for contemporary organizations facing technological disruption, cyber-threats, and systemic interdependence. Yet despite its prominence, organizational resilience remains conceptually fragmented, methodologically underdeveloped, and frequently enacted symbolically rather than operationally. This paper advances a research agenda for organizational resilience in digitally mediated environments, synthesizing insights from resilience theory and information systems scholarship. We argue that as organizational processes become embedded in digital infrastructures and ecosystems, resilience acquires a distinct sociotechnical specification. Identifying three critical gaps—conceptual ambiguity, empirical fragmentation across levels, and limited operationalization—we propose an enactment–scope framework that clarifies how resilience becomes enacted through observable mechanisms and distributed across organizational levels. Cyber resilience is examined as a high-pressure instantiation that renders the digitally structured nature of organizational resilience particularly visible. By anchoring resilience within the dynamics of digital infrastructures and governance arrangements, this paper positions information systems research at the forefront of resilience scholarship.

\end{abstract}

% Keywords
\vspace{15pt}
\noindent \textbf{Keywords:} Resilience, Socio-technical Systems, Information Systems (IS), Organizational Adaptation, Modeling and Simulation, Digital Environments.

\newpage

% Main sections
\section{Introduction}
\label{sec:Introduction}

In early 2023, engineers at Samsung inadvertently exposed proprietary source code and internal documents by entering them into publicly accessible generative artificial intelligence (AI) tools, prompting the company to ban the use of ChatGPT and similar systems among employees to prevent further leaks \autocite{ray2023samsung}. Security firms now report hundreds of generative-AI-related data violations each month across large organizations as employees routinely engage with AI tools outside approved workflows, risking unintentional disclosure of sensitive information and intellectual property \autocite{arora2025understanding}. At the same time, authoritative cybersecurity assessments warn that advanced AI models can automate and accelerate cyberattacks, lowering the barrier to entry for sophisticated threats and enabling malicious actors to scale operations in ways previously constrained by human and technical resources \autocite{arif2024overview}.\\

These developments illustrate a new class of organizational risk: AI-enabled exposures and attack vectors that propagate faster than traditional defenses can detect or govern, and that bypass established security and compliance controls altogether. AI is not alone. Similar dynamics arise with quantum computing, which threatens the durability of cryptographic infrastructures \autocite{lindsay2020demystifying}; internet of things (IoT) and operational technology convergence, which enables cascading cyber-physical failures \autocite{fu2025systematic}; and software supply-chain platforms, where low-cost, scalable exploits can rapidly propagate across ecosystems \autocite{shen2024understanding}.\\

Against this backdrop, resilience has become the digital regulator’s priority \autocite{dupont2019cyberresilience}. As crises such as COVID-19, global cyberattacks, and systemic information technology (IT) failures have exposed the vulnerabilities of digital infrastructures, resilience has shifted from an abstract aspiration to a core regulatory mandate \autocite{almeida2020challenges}. From financial services to healthcare to critical infrastructure, organizations are now expected to anticipate, absorb, and adapt to disruption -- not only to ensure operational continuity but also to maintain legitimacy in a hyperconnected world \autocite{dupont2019cyberresilience, ivanov2022viable}. The growing uncertainty of the cyber landscape, further amplified by rapid AI adoption and the digital acceleration triggered by the pandemic, has fundamentally redefined expectations for digital systems and those who govern them \autocite{he2022building}.\\

A pronounced mismatch persists between regulatory ambition and organizational capability: most organizations struggle to define what resilience means in a digital context, let alone how to operationalize or measure it \autocite{bjorck2015cyber}. While the reality for organizations is increasingly of continuous, novel and overlapping threats and disruptions, organizational structures constraint them in a way that they can only afford to prepare thoroughly scenarii for one or a few exceptional shocks. One root cause lies in static protocols, rigid control frameworks, and audit-driven compliance \autocite{power1997audit, woods2015four}. These practices -- designed for episodic crises -- provide the appearance of control but do little to sustain adaptive performance in environments of persistent uncertainty, an illusion reinforced by audit culture and decision-making under uncertainty \autocite{power1997audit, satinover2007illusion}. Yet, there is a deep sense of dissonance in the mind of practitioners that their organization may not be capable of preparing fast enough for a greatly accelerated future of technological development and dangers that come along \autocite{power1997audit, hollnagel2013resilience}.\\

This discomfort lies at the heart of the resilience implementation gap. Although resilience is widely invoked in standards, risk frameworks, and regulatory guidelines, it is rarely enacted as a lived organizational capability. Organizations continue to default to compliance-based approaches that signal resilience to external audiences while leaving internal practices unchanged \autocite{holling1996engineering, power1997audit, walker2004resilience}. Empirical evidence confirms this pattern. To get a better view of this paradox in practice, we conducted a ligthweight investigation with 18 cybersecurity and risk professionals. 
%While regulatory expansion through instruments such as DORA, NIS2, the Cyber Resilience Act, or FINMA’s Circular 2023/1 has become the primary driver of initiatives, these efforts often translate into audit preparation and budget compliance rather than into adaptive routines, cross-departmental learning, or proactive culture change.\\  
%To illustrate these tensions, Table \ref{tab:RotaResult} synthesizes recurrent themes emerging from both the literature and our field evidence. 
Table \ref{tab:RotaResult} highlights six patterns: (i) the persistence of ambiguities in the definition of resilience and further in how resilience could be quantified, (ii) the dominance of regulatory pressure and symbolic compliance \autocite{power1997audit}, (iii) the endurance of siloed organizational structures, (iv) the underutilization of learning from incidents, (v) the uneven role of human and cultural factors, and (vi) the paradox between the regulatory injunction for flexibility and organizational demand for precision (see Appendix \ref{appendix:C} for further results and methodological details). Together, these patterns capture the practical contradictions that prevent resilience from being enacted as an organizational capability. Rather than treating resilience as a singular construct, they reveal it as a fragmented phenomenon, somewhat contested within the organization, and surely invoked widely but inconsistently translated into practice.\\


\begin{table}[H]
    \centering
    \singlespacing
    \begin{tabular}{p{0.25\textwidth} p{0.35\textwidth} p{0.35\textwidth}}
        \toprule
        \makecell[tl]{\textbf{Theme from Field} \\  \textbf{Evidence}} & \textbf{Dimension of Resilience in Organization}  &  \textbf{Implications for Practice}\\
        \midrule
        \makecell[tl]{\textbf{Quantification \&} \\ \textbf{Ambiguity}} & \makecell[tl]{Organizations lack operational\\ metrics; resilience remains rhe- \\torical and fragmented across \\ units.} & \makecell[l]{Shared definitions and standar- \\ dized constructs are needed to \\ enable measurement.}\\
         \makecell[tl]{\textbf{Regulatory Pressure} \\ \textbf{\&} \\ \textbf{Symbolic Compliance}} &
        \makecell[tl]{Regulatory audits dominate, \\ producing “tick-box” resilience \\ without adaptive capacity.} &
        \makecell[tl]{Tools must assess adaptive be- \\havior, not just documentation.} \\
        \makecell[tl]{\textbf{Coordination \&} \\ \textbf{Silos}} & \makecell[tl]{IT, risk, compliance units oper- \\ate in silos; limited cross-func- \\tional coordination.} & \makecell[l]{Organizational design should \\ support cross-unit collaboration \\ under crisis.}\\
        \makecell[tl]{\textbf{Learning from} \\ \textbf{Incidents}} & \makecell[tl]{Lessons learned rarely institu-\\tionalized; ad-hoc reviews lack \\ persistence.} & \makecell[l]{Mechanisms are needed to embed \\ feedback and retain knowledge.}\\[3em]
        \makecell[tl]{\textbf{Human \&} \\ \textbf{Cultural Factors}} & \makecell[tl]{Expertise, simulation exercises,\\ and culture matter, but interpre- \\tations vary across units.} & \makecell[l]{Behavioral and cultural aspects \\should be integrated into resilience\\ frameworks.}\\
        \makecell[tl]{\textbf{Flexibility–Precision} \\ \textbf{Paradox}} & \makecell[tl]{Regulators provide flexible defi-\\nitions; firms request precise \\ ones.} & \makecell[l]{Balance between contextual adap-\\tation and measurable guidance \\ must be achieved.}\\
         \bottomrule
    \end{tabular}
    \caption{Field evidence on how organizations interpret and enact resilience in practice (Method : semi-structured interviews, c.f. Appendix \ref{appendix:C} for further results and methodological aspects).}
    \label{tab:RotaResult}
\end{table}

The challenge is compounded by conceptual ambiguity. As resilience has diffused across disciplines and industries, its meaning has fragmented -- ranging from system up-time to business continuity, regulatory posture, or security architecture. This broadening definition further undermines coherence and accountability \autocite{duchek2020organizational, woods2015four}. Even when frameworks exist, incidents and near-misses are seldom leveraged for organizational learning, leaving reflexivity, decentralization, and coordination underdeveloped \autocite{barasa2018what, francis2014metric}. Standards and guidelines continue to proliferate \autocite{linkov2019fundamental}, but organization-ready methods to quantify, simulate, and design resilience remain missing \autocite{kuypers2018designing}. In practice, this means resilience risks remaining an attractive principle with little actual traction on the ground.\\
\\

These insights reinforce the central motivation of this work: resilience is not absent from organizational discourse but enacted in ways that are partial, instrumental, and misaligned with its systemic intent. Efforts remain largely driven by external pressures rather than by endogenous capacity building, leaving organizations without a clear, actionable definition of what it means to build resilience.\\ 
\\

These insights reinforce the central motivation of this work: resilience is not absent from organizational discourse but enacted in ways that are partial, instrumental, and misaligned with its systemic intent. Efforts remain largely driven by external pressures rather than by endogenous capacity building, leaving organizations without a clear, actionable understanding of what it means to build resilience. At its core, resilience requires organizations to operate under uncertainty, adapt continuously, and accept limits to centralized control -- conditions that are organizationally costly, culturally resisted, and therefore difficult to institutionalize. A more robust perspective therefore situates resilience as a socio-technical capability: distributed across actors, sustained through coordination and trust, and enabled by adaptive governance arrangements rather than static control mechanisms.\\
\\

To address this impasse, the paper adopts a deliberately integrative and problem-driven approach. Rather than offering yet another definition of resilience or prescribing best practices, we examine why resilience repeatedly collapses into compliance, symbolism, or fragmentation once translated into organizational settings. We treat ambiguity and paradox not as conceptual failures to be eliminated upfront, but as empirical signals of deeper socio-technical tensions between regulation and practice, control and adaptation, and measurement and learning. By systematically tracing how resilience has been conceptualized, operationalized, and enacted across disciplines -- and where these translations break down in practice -- we shift the focus from what resilience should mean to how it can be rendered observable, comparable, and cumulative without undermining its adaptive nature.\\
\\

Our contribution is to establish resilience as a measurable, socio-technical capability central to the IS field. We move beyond dominant treatments of resilience as rhetoric or regulatory compliance by synthesizing cross-disciplinary insights and organizational evidence to show why resilience consistently fails to materialize in practice. We identify the core conceptual, methodological, and empirical gaps that impede operationalization, propose an integrative framework grounded in behavioral indicators, simulation, and institutional dynamics, and advance a research agenda that transforms resilience from an aspirational label into a tractable theoretical and empirical construct. In doing so, this paper positions resilience as a key frontier for IS scholarship on organizational adaptation under persistent digital disruption.\\
\\

The paper proceeds as follows. Section \ref{sec:Background} synthesizes how resilience has been conceptualized across disciplines and within management and information system research, highlighting how its diffusion has generated persistent conceptual ambiguity. Section \ref{sec:RG} identifies the resulting conceptual, methodological, and empirical gaps that explain why resilience is widely proclaimed yet rarely enacted in practice. Section \ref{sec:RD} advances a research agenda aimed at overcoming this impasse by developing measurable indicators, simulation-based methods, and integrative frameworks that render resilience observable and tractable. Section \ref{sec:ImplicationsResearch} discusses the implications of this agenda for both IS scholarship and organizational practice, emphasizing how resilience can be cultivated as a genuine capability rather than displayed as a compliance artifact. Section \ref{sec:Conclusion} concludes by restating the contribution of the review and outlining pathways toward cumulative and actionable resilience research in the IS field.\\




\section{\MakeUppercase {Background}}
\label{sec:Background}
\renewcommand{\baselinestretch}{2.0}\normalsize

This section examines how resilience has been conceptualized across disciplines and how those ideas were later taken up in management and information systems research. Resilience did not travel from one field to another as a stable concept. It developed through distinct intellectual traditions, each carrying its own assumptions about stability, adaptation, control, and system change. Those differences matter in organizational settings because they shape what resilience is taken to mean, how it is measured, and what kinds of interventions are treated as legitimate or effective. Tracing these differences helps explain why resilience has become so widely invoked in organizations while remaining difficult to define, compare, and operationalize consistently.

\subsection{How Resilience Was Shaped Across Disciplines}

\paragraph{Mathematics and Physics}

The earliest foundations of resilience emerged in mathematics and physics, where it was generally understood as the capacity of a system to absorb disturbance and return to equilibrium. Early formulations in physics and systems thinking emphasized elasticity, stability, and homeostatic feedback as defining features of resilient behavior \autocite{wiener1948cybernetics,panteli2017metrics,rankine1858manual}. In that view, resilience is mainly about rebound and temporal recovery. This perspective proved especially useful for formal and quantitative treatments of engineered systems because disturbance-response dynamics can be specified with considerable precision. Its limitation is equally clear: it leaves relatively little room for cases in which systems do not return to a prior stable state.

\paragraph{Biology and Ecology}

Biology and ecology changed the meaning of resilience more fundamentally. The emphasis moved away from equilibrium restoration and toward persistence through variability, adaptation, and regime change. Resilience no longer meant return to a prior state so much as the capacity to reorganize while preserving core functions \autocite{holling1973resilience,gunderson2002panarchy}. This perspective stresses that resilience emerges through feedbacks across ecological, social, and institutional levels \autocite{sundstrom2014transdisciplinary}. Ecological models such as the adaptive cycle illustrate transitions between stability regimes and show how diversity and redundancy can buffer systems against collapse \autocite{walker2004resilience,folke2004regime,folke2010resilience}. Early metrics such as return variance, autocorrelation, and recovery time were used to characterize resilience under uncertainty \autocite{ives1995measuring} (see Appendix \ref{appendix:A2}). This tradition matters because it treats resilience as persistence through reorganization rather than simple recovery. At the same time, it creates a problem for later organizational research: once resilience includes transformation as well as restoration, the relevant baseline, success criteria, and performance measures become harder to specify.\\

These perspectives are often treated as compatible, but they are not built on the same assumptions. Engineering resilience is oriented toward recovery to a known and valued state, whereas ecological resilience accepts irreversibility, nonlinear change, and the possibility that persistence may require transformation rather than restoration \autocite{holling1996engineering,walker2004resilience}. A further extension of this trajectory, articulated by \textcite{taleb2012antifragile}, treats some systems as not merely robust to disturbance but \emph{antifragile}: capable of gaining structurally from variability, exposure, and stress. In organizational settings the antifragile case highlights mechanisms—deliberate exposure to controlled stress, redundancy designed for option value, learning loops that compound across incidents—that can actively strengthen adaptive capacity over time, complementing both restoration- and transformation-oriented accounts. That spectrum becomes difficult to ignore in organizational research, where resilience is often invoked as if stability, adaptation, transformation, and antifragility formed a single coherent objective \autocite{woods2015four,duchek2019organizational}.

\paragraph{Neuroscience and Psychology}

Research in neuroscience and psychology extends the concept again by framing resilience as the capacity of individuals and groups to endure, adapt, and sometimes grow through adversity \autocite{luthar2000construct,kalisch2024neurobiology}. This work emphasizes mechanisms such as cognitive flexibility, learning, and emotional regulation as central to adaptive responses under stress \autocite{masten2010disasters,feder2019biology}. Empirical studies further connect resilience to stress-regulation processes and to the plasticity of adaptive capacities shaped by experience over time \autocite{mcewen2004protection,rakesh2019resilience}. Psychometric instruments such as the Wagnild and Young Scale \autocite{wagnild1993development} and the Connor-Davidson Resilience Scale \autocite{campbell2007psychometric} operationalize resilience mainly at the individual and group levels. Moving beyond strictly individual accounts, \textcite{ungar2021modeling} treats resilience as a multisystemic phenomenon distributed across personal, social, institutional, and ecological contexts. In organizational settings, especially in domains such as cybersecurity operations, healthcare, and emergency response, these insights help explain how stress, judgment, coordination, and learning shape adaptive action under pressure \autocite{norris2007community,kruk2017building,haldane2021health}. What remains less developed is the link to socio-technical organization research. Individual resilience has been studied extensively, but far less is known about how these adaptive capacities scale to collective organizational processes or interact with technological infrastructures.

\paragraph{Institutional Economics}

A major exception to individual-centered accounts appears in institutional economics, which places the collective and governance dimensions of resilience at the center. In her foundational analysis of common-pool resources, \textcite{ostrom1990governing} showed how communities operating under uncertain and demanding conditions can sustain shared resources through self-organized governance arrangements grounded in trust, reciprocity, monitoring, and graduated sanctions. In this tradition, resilience is not merely recovery or endurance. It is the capacity of a socio-economic system to sustain cooperation under uncertainty, external pressure, and resource constraints. Research on social-ecological systems similarly shows that long-term resilience depends on adaptive rule-making, polycentric governance, and alignment between incentives, actors, and shared resource systems \autocite{ostrom1990governing,lebel2006governance}. This tradition is important because it treats resilience as an emergent outcome of institutional design and collective action rather than as an individual trait or a purely technical property. What it leaves more open is how these governance principles should be translated into contemporary organizations whose coordination increasingly depends on digital infrastructures and interorganizational dependencies.

\paragraph{Cybernetics and Engineering}

Control theory and engineering reintroduced resilience into applied domains such as infrastructure, safety, and risk management. Here resilience was typically defined as the ability of a system to recover functionality after disruption, often measured through recovery time, reliability, and functional uptime \autocite{bruneau2003framework}. A foundational engineering contribution distinguishes resilience as recovery speed, reliability as frequency of success, and vulnerability as severity of failure \autocite{hashimoto1982reliability}. This equilibrium-oriented view remains influential where rapid restoration of critical functions is the primary objective (see Appendix \ref{appendix:A1}). In resilience engineering, \textcite{paries2017resilience} identified four practical capabilities -- responding, monitoring, learning, and anticipating -- as key levers for improving system resilience. This tradition offers strong operational clarity and a clear evaluative logic. Its limits become more visible, however, in settings marked by adversarial adaptation, non-stationary environments, and strategic contestation, such as cybersecurity and digitally networked systems \autocite{dupont2019cyberresilience}.

\paragraph{Complexity and Socio-technical Systems}

Complexity science and socio-technical systems theory shift the focus once more by emphasizing emergence, self-organization, and nonlinear interaction. In this view, resilience arises less from centralized control than from local interaction among heterogeneous actors and from feedback mechanisms operating across scales \autocite{levin1998ecosystems,krakovska2023resilience}. Modeling traditions such as agent-based simulation, network analysis, and game-theoretic approaches help explain how system-level resilience can emerge from distributed behavior \autocite{albert2000error,myerson2013game,gao2016universal}. Research on network robustness, for example, shows that scale-free systems may be resilient to random failures yet highly vulnerable to targeted attacks \autocite{albert2000error}. Cross-scale work on infrastructure further shows that resilience must be understood across spatial, temporal, and governance levels because disruptions frequently cascade across them \autocite{moreno2018womens}. Social-ecological systems research complements this perspective by stressing the role of institutions, collective learning, adaptive governance, and actor-resource-rule configurations in sustaining resilience \autocite{berkes1998linking,anderies2004framework,adger2005socialecological}. The structure of social networks matters as well because it conditions information flows, coordination capacity, and collective action \autocite{bodin2009role}. More recent work also makes clear that resilience strategies entail political, social, and economic trade-offs \autocite{sornette2013exploring,aase2020resilience,leichenko2024climate}. This perspective is especially useful for organizational and digital contexts because it foregrounds distributed adaptation and interdependence. Its drawback is different: it is often easier to identify resilience as an emergent system-level property than to translate that property into clear organizational responsibilities, interventions, or metrics.\\

These traditions do not converge on a single understanding of resilience. Some define it as recovery to equilibrium, others as persistence through reorganization, others as a function of cognition, governance, or distributed interaction, and Taleb's antifragility extends the trajectory further by treating exposure to disorder as a potential source of structural gain rather than as a property to be neutralized \autocite{taleb2012antifragile}. Resilience therefore entered organizational research not as a unified construct but as a composite label carrying partly incompatible assumptions about what should persist, what may change, who adapts, and how adaptation should be evaluated. Those tensions are central to understanding why resilience became attractive in organizational discourse while remaining difficult to define and operationalize consistently.

\subsection{Resilience in Organizations}

Organizational research brought resilience much closer to managerial practice. In this literature, resilience is no longer treated only as a broad system property. It is translated into questions of capability, design, governance, and adaptation. That translation was productive, but also selective. Management and organization scholars drew on several resilience traditions at once without fully resolving their underlying differences, especially over what resilience consists of, how it should be measured, and whether it refers mainly to stability, adaptation, or transformation.

\subsubsection{The Construction of Resilience in Organizations}

A large part of the organizational literature explains resilience through two broad sets of mechanisms: structural arrangements and cultural practices. Together, these strands describe how organizations prepare for disruption, respond under pressure, and learn over time \autocite{bhamra2011resilience,weick2011managing,duchek2019organizational}.\\

Structural arrangements matter because they shape how control, information, and resources move when conditions become unstable. Decentralized structures can support quicker local responses when disruptions unfold unevenly across units or time horizons \autocite{boin2013resilient}. Modularity can limit cascading failure by isolating breakdowns and preserving partial system functionality. Ambidextrous forms of organization can also help firms maintain efficiency while preserving enough flexibility to innovate under pressure \autocite{bhamra2011resilience}. These features are especially relevant in digital and cybersecurity settings, where overly centralized coordination can create bottlenecks and increase vulnerability.\\

Cultural practices matter no less. Rehearsals, simulation-based training, and pre-mortem analysis help build shared mental models, reduce uncertainty, and strengthen trust under stress \autocite{westrum2004typology,klein2008performing,weick2011managing}. Practices such as improvisation, learning orientation, and psychological safety further support teams facing novel or high-pressure situations \autocite{lanzara1999transient,vogus2007organizational}. In fast-moving threat environments, these practices make it possible to react and learn before formal control systems can be revised \autocite{westrum2004typology,duchek2019organizational}.\\

One of the main contributions of this literature is that it makes resilience organizationally interpretable. Resilience is no longer treated only as a system-level outcome but as a property that can be shaped through deliberate design choices, organizational routines, and cultural conditions. The remaining task is to identify, more precisely, the pathways through which these structural and cultural enablers actually produce adaptive behavior, and the indicators that would confirm whether such behavior is taking place.

\subsubsection{Resilience as a Dynamic Managerial Capability}

A second major line of work treats resilience as a dynamic managerial capability. From this perspective, resilience is less a reactive outcome than a path-dependent capacity formed through routines, governance arrangements, resource deployment, and strategic intent \autocite{lengnick2011developing,boin2013resilient}. This marks a shift away from narrow views of robustness or risk avoidance. The emphasis falls instead on an organization's ability to interpret uncertainty, reconfigure resources, and continue functioning under pressure.\\

Recent work explicitly frames organizational resilience as a form of dynamic capability that allows firms to sense disruption, respond effectively, and in some cases derive advantage from it \autocite{teece2007explicating,duchek2019organizational}. In that view, resilience depends on microfoundations such as routines, heuristics, and cognitive schemas that shape how actors interpret events and act under uncertainty \autocite{hepfer2022heterogeneity}. Resilience therefore cannot be reduced to resistance or continuity alone. It requires an adaptive orientation embedded in both operational and strategic action.\\

In practice, that capability is usually supported by combinations of managerial levers. Slack resources and redundancy can provide buffering capacity \autocite{barasa2018what}. Ambidextrous arrangements can support the simultaneous pursuit of stability and innovation \autocite{bhamra2011resilience}. Resilience audits and readiness diagnostics aim to assess whether organizations can absorb and adapt to shocks \autocite{burnard2011organisational,calder2021iso}. This stream is useful because it translates resilience into a language managers can act on. What it still does not resolve is a basic problem: resilience is often described as a desirable capability before there is clear agreement on what observable performance would actually demonstrate that the capability exists.

\subsubsection{Cybersecurity Resilience as a Strategic and Institutional Imperative}

Cybersecurity is a particularly revealing setting because many of the broader organizational tensions around resilience become more visible there. Cyber risks are marked by uncertainty, rapid change, and interdependent vulnerabilities that cut across technical, human, and organizational layers \autocite{linkov2013resilience}. They are also often asymmetric: threats may emerge quickly, evolve during response, and remain only partly visible until damage is already underway. Under such conditions, compliance-oriented and control-oriented approaches are not sufficient on their own \autocite{dupont2019cyberresilience}. This is one reason cybersecurity is now increasingly treated not only as a control problem, but as a resilience problem \autocite{linkov2013resilience,dupont2019cyberresilience}.\\

The cyber-resilience literature generally presents resilience as an integrated capability combining technological safeguards, strategic foresight, organizational adaptability, and institutional alignment \autocite{bjorck2015cyber,lostri2022shared}. Standards such as ISO 22301 emphasize the maintenance of core functions under disruption \autocite{calder2021iso}, while other work places more weight on learning, reflexive governance, and ongoing adjustment \autocite{hausken2020cyber}. Regulatory pressures, including those associated with Basel III and related resilience initiatives in finance, have further reinforced resilience as a strategic imperative by tying it to legitimacy, continuity, and systemic risk management \autocite{baselcommittee2011basel,almeida2020challenges}. This has contributed to a broader shift in governance: firms are increasingly judged not only by the strength of their own security posture, but also by the extent to which they help sustain resilience across wider ecosystems \autocite{bjorck2015cyber,lostri2022shared,gillard2023efficient}.\\

At the same time, this work also clarifies where current organizational thinking can be extended. Formal controls, standards, and regulatory mandates have done important work in clarifying accountability and guiding resource allocation; their next stage of development concerns how compliance and document-centered assurance can be coupled more directly to observable adaptive capacity in practice. Initial evidence from the field, including the interviews summarized in Appendix~\ref{appendix:C}, suggests this coupling is uneven and warrants closer empirical attention. AI, predictive analytics, and scenario modeling offer promising tools to support anticipation and adaptive response \autocite{kuypers2018designing,gillard2023efficient}, and their integration into routine organizational governance is steadily progressing. Across these developments, a productive tension remains between resilience as regulatory requirement, resilience as managerial capability, and resilience as an emergent property of socio-technical coordination under pressure \autocite{boin2013resilient,linkov2013resilience,duchek2019organizational}. Reconciling these three faces of resilience is precisely what the present review aims to support.\\
\\
Together, management and cybersecurity research have made resilience visible in organizational practice by identifying structural, cultural, strategic, and governance-related conditions that support adaptive action. Building on these contributions, three open questions continue to shape the broader resilience literature: how to refine the concept itself, how to converge on what counts as observable resilient performance, and how to strengthen the link between conceptual claims and operational measures. Engaging these questions is what allows resilience, already a serious organizing principle in policy and practice, to be defined and enacted with greater consistency in digitally mediated settings.
% \section{Resilience in Management}

\subsection{Resilience as a Dynamic Managerial Capability}

Among the many domains contributing to resilience thinking, management plays a pivotal role in operationalizing and institutionalizing resilience across complex organizational landscapes. Within this domain, resilience is not merely a theoretical construct but a lived organizational reality that is shaped through practices, routines, and governance arrangements. Management research has gradually shifted from conceptualizing resilience as a reactive outcome to treating it as a proactive, path-dependent capability. This evolution reflects the broader theoretical progression in the field—from classical contingency-based models to more contemporary views rooted in dynamic capabilities, complex systems, and institutional embeddedness.\\ 

Organizational resilience is increasingly framed as a dynamic capability that allows firms to sense, anticipate, respond to, and even leverage disruptions for transformation \autocite{duchek2020organizational, teece2007explicating}. The dynamic capabilities framework emphasizes the processes by which firms purposefully create, extend, or modify their resource base to meet environmental challenges. Within this framing, resilience emerges from microfoundational elements, including routines, decision-making heuristics, and cognitive schemas that shape how actors interpret and respond to uncertainty \autocite{hepfer2022heterogeneity}. It is not simply about robustness or risk mitigation; rather, it is an adaptive orientation embedded in the fabric of strategic and operational activity. The interplay between slack resources, distributed leadership, and organizational identity often underpins this capacity, enabling firms to deploy resources flexibly and sustain performance in the face of adversity \autocite{bhamra2011resilience, barasa2018what}. In practice, this can manifest through investment in absorptive capacity, dual structures that promote both innovation and efficiency, and intentional resilience audits that assess an organization's readiness to absorb and adapt to disruptive shocks.

\subsection{Design, Culture, and the Construction of Resilience}
From an organizational design perspective, resilience is fostered through both structural and cultural mechanisms. Structural enablers include modular architectures, redundancy in critical functions, and decentralization of authority to support rapid decision-making \autocite{lengnick2011developing}. These allow systems to reconfigure rapidly and isolate failures, limiting their propagation across units. Culturally, resilience is nurtured through norms that value learning, reflexivity, and psychological safety. A resilient culture embraces ambiguity, empowers employees to experiment, and promotes distributed sensemaking across hierarchical levels. Leadership plays a critical role in this domain, as resilient organizations often display leadership that is forward-looking, emotionally intelligent, and capable of mobilizing cross-functional alignment \autocite{gillard2023efficient}.\\

Participatory mechanisms—such as hackathons, cross-disciplinary task forces, or innovation labs—further reinforce this adaptive culture by facilitating the emergence of novel solutions and strengthening collective sensemaking \autocite{maillart2024computational}. These configurations also illustrate how resilience is co-constructed through iterative negotiation among actors, practices, and technologies. Over time, such practices may crystallize into organizational routines that serve as repositories of institutional memory and adaptive learning. Together, these structural and cultural elements underscore that resilience is not exogenous or incidental but instead deeply embedded in how an organization configures its resources, cultivates its people, and governs its practices.

\subsection{Cybersecurity Resilience as a Strategic and Institutional Imperative}

The cybersecurity domain serves as a particularly illustrative context in which these resilience dynamics unfold. Cyber risks are characterized by high uncertainty, rapid evolution, and interdependent vulnerabilities that cut across technical, human, and organizational layers \autocite{linkov2013resilience}. Unlike traditional risk scenarios that are probabilistic and often bounded in scope, cybersecurity threats are emergent, asymmetrical, and often invisible until activated. Traditional risk management approaches, focused on compliance and control, are often inadequate in the face of such complexity. Consequently, cyber resilience is now viewed as an integrated capability encompassing not only technological safeguards but also strategic foresight, organizational agility, and institutional alignment \autocite{bjorck2015cyber, lostri2022shared}.\\

Frameworks such as ISO 22301 provide structured guidance for crisis response, emphasizing the importance of maintaining core functions under stress \autocite{iso201922301}. However, emergent views stress the need for adaptive processes rooted in learning and reflexive governance \autocite{hausken2020cyber}, where security practices evolve in tandem with threat landscapes. Regulatory pressures—such as those stemming from Basel III—have institutionalized resilience as a strategic imperative, linking it to organizational legitimacy and systemic risk management \autocite{baselcommittee2011basel, almeida2020challenges}. This has catalyzed a broader governance turn in cybersecurity management, where firms are increasingly accountable not only for their individual risk posture but also for their contribution to systemic resilience.\\

Meanwhile, advancements in AI, predictive analytics, and scenario modeling empower firms to anticipate cyber threats and orchestrate adaptive responses in real time \autocite{kuypers2018designing, gillard2023efficient}. These technologies allow for the simulation of threat vectors, visualization of attack surfaces, and real-time adjustment of defensive protocols. In sum, cyber resilience exemplifies the synthesis of technological infrastructure, human cognition, and institutional design --demonstrating how resilience in management is both a means of navigating uncertainty and a strategy for sustained adaptation. When viewed through the lens of resilience, cybersecurity becomes less about perimeter defense and more about the dynamic coordination of resources, people, and information to enable anticipatory and responsive action at scale.
\section{Research Gaps}
\label{sec:RG}

Despite its prominence in policy, strategy, and academic discourse, resilience remains inconsistently operationalized and unevenly assessed in organizational settings. This challenge does not stem from a lack of theoretical development. As shown above, resilience has been extensively conceptualized and modeled across disciplines, including engineering, ecology, psychology, and complex systems. However, these contributions rest on distinct ontological assumptions and methodological traditions that rarely converge into a unified organizational framework. As resilience diffuses into managerial and regulatory domains, this plurality becomes consequential: organizations are expected to enact resilience without clear guidance on which assumptions, metrics, or governance logics should prevail.\\

As shown above, resilience has been used to denote recovery, robustness, adaptation, learning, or transformation, often simultaneously and within the same organizational setting. These meanings rest on incompatible assumptions about system behavior, temporality, and control: recovery- and robustness-oriented views presuppose a stable baseline to which systems should return, whereas adaptive and transformative views assume non-stationarity, continual change, and the possibility that post-disruption states differ fundamentally from prior ones. When imported into organizations, this conceptual plurality can produce misalignment between investment, evaluation, and governance logics: organizations may invest in robustness while being evaluated on adaptability, or implement compliance controls while being rhetorically charged with transformation.\\

Our review identifies three interlinked gaps that emerge not from theoretical absence, but from insufficient integration and translation into organizational research and practice: conceptual plurality without operational convergence, methodological fragmentation across domains, and limited empirical validation in organizational contexts. {\it Conceptually}, the absence of an integrated operational core within management and information systems research prevents cumulative theory-building. {\it Methodologically}, existing metrics and formal models are rarely embedded in organization-level governance and evaluation systems, and therefore struggle to capture resilience as a dynamic, socio-technical capability. {\it Empirically}, the lack of standardized and longitudinal indicators limits comparability and learning across cases. Among these challenges, the development of operationalizable and transferable measures remains foundational \autocite{woods2015four}, a concern repeatedly documented across domains ranging from infrastructure to healthcare \autocite{heinimann2017infrastructure, kruk2017building}. Details of the systematic review process, covering search, screening, and coding across 340 publications, are provided in Appendix~\ref{appendix:B}, while illustrative field evidence of how these tensions manifest in organizational practice is synthesized in Appendix~\ref{appendix:C}.\\
\\
First, {\bf conceptual plurality without operational convergence:} Resilience has been conceptualized across domains through lenses of mechanical stability, ecological transformation, cognitive adaptability, and structural flexibility (see Appendix \ref{appendix:A}). While each view offers valuable and often rigorously formalized insights, their assumptions are not equivalent. Engineering resilience emphasizes restoration of predefined functions; ecological resilience foregrounds regime shifts and adaptive reorganization; psychological resilience centers on coping and growth under stress; socio-technical perspectives highlight emergence and distributed coordination. Although these perspectives can be interpreted as complementary at an abstract level, their implications for organizational design, measurement, and accountability differ substantially.\\

In practice, this plurality can generate definitional drift. Terms such as resistance, robustness, recovery, agility, and adaptation are often used interchangeably, despite referring to distinct mechanisms and temporal logics. Consequently, resilience may be invoked symbolically, through audits, strategy documents, or compliance templates, rather than embedded as a coherent adaptive capability \autocite{duchek2020organizational}. Without clearer articulation of how these perspectives translate into measurable organizational mechanisms, resilience risks remaining conceptually rich yet operationally diffuse.\\
\\
Second, {\bf methodological fragmentation and limited transfer:} This gap is visible in the cross-disciplinary classification of 89 resilience-related publications (see Table \ref{tab:ResearchGaps}). Simulation models, network analyses, and quantification frameworks are well established in engineering, complex systems, and climate research \autocite{francis2014metric}. However, their systematic integration into domains central to organizational and managerial resilience -- such as cybersecurity, risk management, and strategic governance -- remains limited. Table \ref{tab:ResearchGaps} shows that \textit{Markers \& Metrics} and \textit{Technology} are comparatively underrepresented in these fields, and that simulation-based approaches are rarely incorporated into organization-level governance frameworks \autocite{linkov2013resilience, kuypers2018designing}. The result is a continued reliance on maturity models, expert judgement, or audit documentation—approaches that may support compliance but rarely capture adaptive performance under real-world conditions.\\

Even within cybersecurity, formal models of attack–defense dynamics and adaptive response strategies exist \autocite{aydin2018integration}, but are seldom integrated into broader organizational resilience architectures. Other transferable approaches, such as ecological network analysis that quantifies interdependencies \autocite{luke2012systems}, transportation frameworks combining metrics and recovery modeling \autocite{jacobs2017applying}, or agent-based methods for infrastructure risk assessment \autocite{dawson2011agentbased}, remain under-exploited in digital resilience planning. The challenge is therefore not methodological scarcity, but limited cross-domain integration: technical innovations in measurement are well developed but insufficiently translated into organization-ready governance tools \autocite{heinimann2017infrastructure, aerts2018integrating}.\\
\\
\begin{center}
\newgeometry{top=.15cm, bottom=2cm, left=2.3cm, right=2.3cm}
\singlespacing
\begin{landscape}
\tiny
\begin{longtblr}[
  caption = {\textbf{Classification of Resilience-related Publications.} This table provides the classification of 89 relevant, impactful and supporting publications associated with resilience according to their corresponding topic/-s and used methodology/-ies. This table allows to extract the key research gaps, where no publication was found (red cells) or only very few (orange cell). Topic-methodology intersections in grey show other research that are less relevant for this research.},
  label = {tab:ResearchGaps}
]{
  width=\linewidth,
  colspec={|Q[1.5cm] Q[1.7cm]|Q[2.6cm] Q[2.6cm] Q[2.6cm] Q[2.6cm] Q[2.6cm] Q[2.6cm] Q[2.6cm]|},
  rowhead = 2,
  hlines, vlines
}
\SetCell[r=2, c = 2]{l} \vspace{2em}\par \textit{Topics}& & \SetCell[r = 1, c = 7]{halign = l}{\textit{Methodologies}} \\
& & \textbf{\makecell{Simulation \\ \& \\ Modeling}} & \textbf{\makecell{Network \\ Analysis}} & \textbf{\makecell{Markers \\ \& \\ Metrics}} & \textbf{\makecell{Technology}} & \textbf{\makecell{Case Studies}} & \textbf{\makecell{Longitudinal \\ Studies}} & \textbf{\makecell{Surveys}} \\
\hline
% === Example block start ===
\SetCell[r=1, c=2]{l}{\textbf{\makecell{Complex Adaptive Systems}}} & &  \SetCell[r=1, c = 1]{l} {\cite{holling1973resilience} \\ \cite{gao2016universal} \\ \cite{pumpuni2017resilience} \\ \cite{krakovska2023resilience}} & \SetCell[r=1, c = 1]{l}{\cite{albert2000error} \\ \cite{bodin2009role} \\ \cite{gao2016universal} \\ \cite{pumpuni2017resilience}} & \SetCell[r=1, c = 1]{l} {\cite{anderies2004framework} \\ \cite{sundstrom2014transdisciplinary} \\ \cite{pumpuni2017resilience}} & \SetCell[r=1, c = 1]{bg = gray!50} & \SetCell[r=1, c = 1]{l} {\cite{ostrom1990governing} \\ \cite{lanzara1999transient} \\ \cite{folke2004regime} \\ \cite{weick2011managing}} & \SetCell[r=1, c = 1]{bg = gray!50} & \SetCell[r=1, c = 1]{l}{\cite{ostrom1990governing} \\ \cite{levin1998ecosystems} \\ \cite{gunderson2002panarchy} \\ \cite{walker2004resilience} \\ \cite{heinimann2017infrastructure}}\\
\hline
\SetCell[r=1,c=2]{l}{\textbf{\makecell{Engineering}}} & & \SetCell[r=1, c = 1]{l}{\cite{wiener1948cybernetics} \\ \cite{panteli2017metrics}} & \SetCell[r=1, c = 1]{l}{\cite{albert2000error} \\ \cite{gao2016universal}} & \SetCell[r=1, c = 1]{l}{\cite{bruneau2003framework} \\ \cite{francis2014metric} \\ \cite{panteli2017metrics}} & \SetCell[r=1, c = 1]{l}{\cite{rane2024artificial}} & \SetCell[r=1, c = 1]{l}{\cite{paries2017resilience}} & \SetCell[r=1, c = 1]{bg = gray!50} & \SetCell[r=1, c = 1]{l}{\cite{holling1996engineering} \\ \cite{francis2014metric} \\ \cite{ivanov2022viable}} \\
\hline
\SetCell[r=1,c=2]{l}{\textbf{\makecell{Ecology \& Social Systems}}} & & \SetCell[r=1, c = 1]{l}{\cite{holling1973resilience} \\  \cite{hashimoto1982reliability} \\ \cite{ives1995measuring}} & \SetCell[r=1, c = 1]{l}{\cite{bodin2009role}} & \SetCell[r=1, c = 1]{l}{\cite{ives1995measuring} \\ \cite{sundstrom2014transdisciplinary}} & \SetCell[r=1, c = 1]{l}{\cite{maillart2024computational}} & \SetCell[r=1, c = 1]{l}{\cite{berkes1998linking} \\ \cite{lebel2006governance} \\ \cite{mahajan2022participatory}} & \SetCell[r=1, c = 1]{l}{\cite{maillart2024computational}} & \SetCell[r=1, c = 1]{l}{\cite{folke2010resilience} \\ \cite{leslie2013response}\\ \cite{francis2014metric} \\ \cite{allen2018quantifying}}\\
\hline
\SetCell[r=1,c=2]{l}{\textbf{\makecell{Climate Change \& Disaster Management}}} & & \SetCell[r=1, c = 1]{l}{\cite{dawson2011agentbased} \\ \cite{panteli2017metrics} \\ \cite{aydin2018integration}} & \SetCell[r=1, c = 1]{l}{\cite{jacobs2017applying} \\ \cite{aydin2018integration}} & \SetCell[r=1, c = 1]{l}{\cite{bruneau2003framework}} & \SetCell[r=1, c = 1]{l}{\cite{arnold2004informationsharing} \\ \cite{aerts2018integrating}} & \SetCell[r=1, c = 1]{l}{\cite{dawson2011agentbased} \\ \cite{mahajan2022participatory}} & \SetCell[r=1, c = 1]{l}{\cite{moreno2018womens}} & \SetCell[r=1, c = 1]{l}{\cite{alexander2013resilience} \\ \cite{meerow2016defining} \\ \cite{adger2005socialecological} \\ \cite{leichenko2024climate}}\\
\hline
\SetCell[r=1,c=2]{l}{\textbf{\makecell{Psychology \& Neuroscience}}} & & \SetCell[r=1, c = 1]{l}{\cite{manikis2019computational} \\ \cite{ungar2021modeling}} & \SetCell[r=1, c = 1]{bg = gray!50} & \SetCell[r=1, c = 1]{l}{\cite{mcewen2004protection} \\ \cite{rakesh2019resilience}} & \SetCell[r=1, c = 1]{l}{\cite{klein2008performing} \\ \cite{bada2019cyber} \\\cite{kalisch2024neurobiology}} & \SetCell[r=1, c = 1]{l}{\cite{bada2019cyber}} & \SetCell[r=1, c = 1]{l}{\cite{rutter1987psychosocial} \\ \cite{masten2010disasters}} & \SetCell[r=1, c = 1]{l}{\cite{wagnild1993development} \\ \cite{luthar2000construct} \\ \cite{campbell2007psychometric} \\ \cite{feder2019biology}}\\
\hline
\SetCell[r=1,c=2]{l}{\textbf{\makecell{Public Health}}} & & \SetCell[r=1, c = 1]{l}{\cite{brauer2008compartmental} \\ \cite{luke2012systems}}  & \SetCell[r=1, c = 1]{l}{\cite{luke2012systems}} & \SetCell[r=1, c = 1]{l}{\cite{cutter2010disaster} \\ \cite{kruk2017building}} & \SetCell[r=1, c = 1]{l}{\cite{arnold2004informationsharing}} & \SetCell[r=1, c = 1]{l}{\cite{barasa2018what} \\ \cite{haldane2021health}} & \SetCell[r=1, c = 1]{l}{\cite{gibbs2013beyond} \\ \cite{aase2020resilience}} & \SetCell[r=1, c = 1]{l}{\cite{norris2007community} \\ \cite{barasa2018what}}\\
\hline
\SetCell[r=2, c=1]{l}{\vspace{-16.25em}\par \textbf{Management}} & \SetCell[r=1, c=1]{l}{\textbf{Organizational}} & \SetCell[r=1, c = 1]{bg = red!20} & \SetCell[r=1, c = 1]{bg = red!20} & \SetCell[r=1, c = 1]{l}{\cite{kuypers2018designing}} & \SetCell[r=1, c = 1]{l}{\cite{klein2008performing} \\ \cite{bada2019cyber} \\ \cite{almeida2020challenges}} & \SetCell[r=1, c = 1]{l}{\cite{lanzara1999transient} \\ \cite{westrum2004typology} \\ \cite{weick2011managing} \\ \cite{henfridsson2013generative} \\ \cite{barasa2018what} \\ \cite{bada2019cyber} \\\cite{boh2023building}} & \SetCell[r=1, c = 1]{bg = gray!50} &  \SetCell[r=1, c = 1]{l}{ \cite{westrum2004typology}\\ \cite{teece2007explicating} \\ \cite{vogus2007organizational} \\ \cite{bhamra2011resilience} \\ \cite{burnard2011organisational} \\ \cite{lengnick2011developing} \\ \cite{boin2013resilient} \\ \cite{henfridsson2013generative} \\ \cite{skopik2016problem} \\ \cite{dupont2019cyberresilience} \\ \cite{barasa2018what} \\ \cite{duchek2019organizational} \\ \cite{calder2021iso} \\\cite{hepfer2022heterogeneity} \\ \cite{he2022building}}\\
\cline[dashed]{2-9}
&  \SetCell[r=1, c=1]{l}{\textbf{Risk \& Cyber}} & \SetCell[r=1, c = 1]{bg = red!20} & \SetCell[r=1, c = 1]{bg = orange!20}{\cite{gillard2023efficient}} & \SetCell[r=1, c = 1]{l}{\cite{linkov2013resilience} \\ \cite{panteli2017metrics}} & \SetCell[r=1, c = 1]{l}{\cite{aerts2018integrating} \\ \cite{almeida2020challenges}} & \SetCell[r=1, c = 1]{l}{\cite{lostri2022shared} \\ \cite{boh2023building}} & \SetCell[r=1, c = 1]{l}{\cite{kuypers2018designing}} & \SetCell[r=1, c = 1]{l}{\cite{baselcommittee2011basel} \\ \cite{bjorck2015cyber} \\ \cite{woods2015four} \\ \cite{skopik2016problem} \\ \cite{linkov2019fundamental} \\ \cite{dupont2019cyberresilience} \\  \cite{hausken2020cyber}  \\ \cite{calder2021iso} \\ \cite{lostri2022shared}} \\
\end{longtblr}
\end{landscape}
\end{center}
\restoregeometry



Third, {\bf limited empirical validation and longitudinal grounding:} Field data from cybersecurity and risk practitioners indicates that resilience is typically assessed through compliance artifacts -- such as audits and certifications -- rather than by systematically analyzing adaptive behavior during disruption. Advanced tools such as predictive simulations or anomaly detection are piloted in innovation units but often remain decoupled from regulatory evaluation structures \autocite{skopik2016problem}. Informal practices -- parallel communication channels, crisis improvisation, or decentralized escalation -- remain critical to adaptive performance but largely invisible in existing assessments.\\

Similar dynamics are evident in healthcare and public health, where resilience is observed in practice yet only partially institutionalized in measurement systems \autocite{kruk2017building}. Moreover, longitudinal designs that track how resilience evolves across repeated disruptions remain rare. Few frameworks explicitly address the uncertainty, reversibility, and trade-offs inherent in resilience assessment \autocite{allen2018quantifying}. Disciplinary fragmentation further limits cross-fertilization, as methods from control theory or complex systems are seldom transferred into organizational research, and widely used practices such as awareness campaigns often fail to foster measurable behavioral adaptation \autocite{bada2019cyber}.\\
\\
Together, these gaps indicate that resilience research possesses substantial theoretical and methodological resources, yet lacks sufficient integration to support cumulative, organization-ready application. Without standardized, validated, and scalable metrics embedded in governance structures, resilience risks remaining a normative aspiration rather than a systematically cultivated capability. Mathematical approaches to dynamical systems—through basin stability, return dynamics, and perturbation response—provide promising foundations that could be adapted to trace organizational adaptation over time \autocite{krakovska2024resilience}. This is not solely a technical concern. For decision-makers, limited metric integration complicates investment prioritization, performance comparison, and accountability to boards and regulators. For scholars, it constrains intervention testing, replication, and theory-to-practice translation. Advancing resilience research therefore requires not additional definitions, but clearer integration of conceptual foundations, methodological tools, and empirical validation strategies. The next section addresses these gaps directly by clarifying operational mechanisms, advancing simulation-based measurement approaches, and outlining empirical pathways for longitudinal validation.

\subsection{Research in Organizational Resilience in Digital Environments}

Organizational resilience has traditionally been conceptualized as the capacity of firms to absorb shocks, adapt to disruption, and sustain core functions over time \autocite{duchek2020organizational,burnard2011organisational}. However, contemporary organizations operate through digitally mediated processes, interconnected infrastructures, and platform-based ecosystems \autocite{hepfer2022heterogeneity, janssen2006network}. As core activities become embedded in information systems, adaptive capacity is increasingly shaped by technological architectures, data flows, and algorithmic coordination mechanisms \autocite{hepfer2022heterogeneity,liu2022network}.\\

We therefore treat \emph{digital organizational resilience} as the contemporary sociotechnical specification of organizational resilience in digitally mediated environments (Duchek, 2020; Hepfer et al., [year]). Rather than constituting a distinct theoretical domain, digital organizational resilience clarifies how organizational resilience is enacted when adaptive capacity is distributed across digital infrastructures, institutional arrangements, and inter-organizational ecosystems (Janssen et al., 2006; Liu et al., 2022). In this view, resilience remains a dynamic capability (Duchek, 2020), yet its realization increasingly depends on digitally embedded coordination mechanisms and infrastructural interdependencies (Hepfer et al., [year]; Boin \& van Eeten, 2013).

Cyber resilience represents a particularly salient instantiation of this specification (Linkov et al., 2013). Adversarial digital threats expose the extent to which organizational resilience depends on information systems, cross-boundary coordination, and infrastructural interdependence (Dupont, 2019; Ibne Hossain et al., 2020). Studying resilience in cyber contexts thus renders visible the mechanisms through which organizational resilience becomes digitally structured and operationalized within interconnected sociotechnical systems (Linkov et al., 2013; Liu et al., 2022).


\section{\MakeUppercase {Research Directions for Digital Organizational Resilience}}
\label{sec:RD}
\renewcommand{\baselinestretch}{2.0}\normalsize


The Enactment--Scope framework provides an important stepping stone for the understanding of resilience, and its value grows further when it is used to generate concrete research questions. This section develops that next step. The focus here is digital organizational resilience, and especially cyber resilience, because digital environments make the core analytical problems particularly visible. In such settings, the substantial work already accomplished by regulators, standard-setting bodies, practitioners, and researchers can be extended by studying resilience as a socio-technical capability that is built, enacted, and tested under conditions of persistent disruption.\\

This shift extends the object of inquiry. When resilience is also examined through observable behavior, infrastructure, and governance arrangements, it becomes possible to evaluate it more directly and to refine the conceptual commitments practitioners and policymakers have already articulated. When it is treated as cumulative rather than static, questions of learning, memory, reinforcement, and erosion move to the center. And in digitally mediated environments, adaptive capacity is shaped not only by managerial judgment and organizational culture, but also by architectures, platforms, institutional settings, and the ways these elements condition what forms of response are possible.\\

The research agenda that follows is organized around two linked questions. The first concerns \emph{enactment}: through what mechanisms does resilience move from stated intention to operational capability? The second concerns \emph{scope}: at what organizational and institutional levels do these mechanisms operate? These questions build on the earlier diagnosis of resilience research. Where organizations have already developed indicators, governance structures, and evaluative benchmarks, the productive next step is to deepen the connection between observed adaptive performance and decision authority. At the same time, resilience in digitally mediated environments is distributed. It is not located in one team, one platform, or one formal rule. It is produced across routines, infrastructures, institutional arrangements, and inter-organizational ecosystems.\\

This is also where information systems research has a distinctive advantage. The field already addresses digital infrastructures, platform governance, socio-technical coordination, and cross-boundary information flows. These are not peripheral concerns here. They are the substance of the problem. The sections below therefore develop three research directions: operationalizing resilience enactment, examining resilience across organizational and institutional scopes, and tracing how resilience evolves over time.

\subsection{Operationalizing Resilience Enactment}

A first priority is to make resilience observable in practice. That is the central task of the enactment dimension of the framework. Organizations have invested significantly in resilience-related strategic language, governance documents, and policy discourse, and the next step is to develop equally robust ways of identifying resilience as it unfolds during disruption. IS research is well placed to support that step because it can connect behavioral indicators, computational traces, and organizational processes in ways that make adaptive capacity empirically visible.\\
\\
Building on the conceptual ground already established, research can now specify how resilience appears in practice and how it can be distinguished from related constructs such as preparedness or adaptability. Comparative work across domains already suggests that resilience emerges through the interaction of routines, infrastructures, and institutional processes \autocite{pumpuni2017resilience,manikis2019computational}. What remains to be developed is a more systematic basis for turning those observations into measures and design principles that organizations can readily apply.\\

Figure \ref{fig:triangle} supports that move by connecting theory, measurement, and governance as the three domains through which resilience becomes enacted capability. Each side of the triangle captures a different mechanism of translation. \emph{Operational indicators} render resilience empirically observable, complementing the conceptual commitments organizations already articulate. \emph{Design principles} allocate decision rights and structure adaptive responses. \emph{Benchmarks and warnings} align incentives while producing actionable signals. The figure matters less as a visual device than as an analytical claim: resilience becomes more readily governable when these three domains are linked rather than developed separately.\\

A converging line of work across several disciplines now develops behavioral and computational indicators that make resilience empirically assessable. Engineering resilience research has long offered formal candidates---recovery speed, reliability, vulnerability, and post-event functionality \autocite{hashimoto1982reliability,bruneau2003framework,francis2014metric,panteli2017metrics}. Public health has produced indicator suites for community and system resilience \autocite{cutter2010disaster,kruk2017building,haldane2021health}. More recent work in cybersecurity and digital risk management has begun to translate these traditions into organizational settings by characterizing investment--exposure dynamics, coordination performance, and incident-response patterns through quantitative and computational analysis \autocite{kuypers2018designing,gillard2023efficient,maillart2024computational,rane2024artificial}. Indicators such as coordination speed, the ability to reconfigure infrastructures, and the variability of response under pressure can complement the strategic commitments organizations have already made with directly measurable counterparts \autocite{mahajan2022participatory}. They also make comparison possible: organizations can benchmark themselves, and researchers can begin to accumulate evidence across cases. Importantly, developing such indicators is not solely a technical exercise. It also rests on a theoretical account of how behaviors at individual and collective levels scale into system-level outcomes \autocite{luthar2000construct,ungar2021modeling,kalisch2024neurobiology}.\\

Digital environments make this task especially urgent. Disruptions unfold quickly, interactions are distributed, and many adaptive signals are embedded in digital traces rather than retrospective narratives. Work in behavioral analytics and AI-assisted modeling therefore offers a useful path forward, particularly where resilience indicators can be calibrated against observable signatures in digital artifacts \autocite{rane2024artificial}.\\

Computational approaches, including agent-based modeling and simulation, are especially useful here. They allow researchers to explore how coordination patterns, incentive structures, and infrastructural configurations shape the capacity of organizations to absorb, reconfigure, and respond to disruption \autocite{arnold2004informationsharing,gillard2023efficient}. Used carefully, such methods do more than illustrate scenarios. They allow theory to be refined against simulated and empirical evidence and help bridge the gap between conceptual language and organization-level practice. In that sense, operationalization is not only about measurement. It is also about making resilience a capability that can be examined, compared, and deliberately improved.

\subsection{Resilience Across Organizational and Institutional Scopes}

A second research direction develops the scope dimension of the Enactment--Scope framework by asking where resilience mechanisms operate and how they are distributed across organizational, infrastructural, institutional, and inter-organizational settings.\\
\\
In digitally mediated environments, adaptive capacity is inherently distributed. It does not sit neatly inside a single team, technology, or governance layer. Governance architectures can enable resilience, but they can also constrain it by creating ambiguity in decision rights, accountability gaps, and escalation failures \autocite{ostrom1990governing,lebel2006governance}. These frictions may remain latent in stable periods, but they become visible under stress. They appear as delayed responses, constrained options, unclear ownership, and coordination breakdowns.\\

Studying these dynamics calls for approaches that connect institutional theory, regulatory ethnography, and actor-network perspectives with empirical work on digital infrastructures and everyday organizational practice. The point is to understand how resilience is enabled or constrained in the lived realities of digitally coordinated work, complementing the formal schemes of control already developed.\\

Table~\ref{fig:resilience_interface_opportunity} introduces an opportunity matrix that maps six research directions for digital organizational resilience---operational indicators, integrative frameworks, comparative and cumulative work, temporal and process markers, normative or value-based analysis, and governance and incentives---across four scopes: \emph{micro-level routines}, \emph{meso-level infrastructures}, \emph{macro-level institutions}, and \emph{inter-organizational ecosystems}. The cell shading indicates where each direction is currently most tractable and where it requires further institutional or methodological investment. The matrix does not replace the enactment perspective; it complements it. Figure~\ref{fig:triangle} clarifies \emph{how} resilience becomes operationalized through indicators, design principles, and evaluative benchmarks; Figure~\ref{fig:enactment_scope} clarifies \emph{where} those mechanisms operate; the opportunity matrix indicates, for each research direction, \emph{where the next round of evidence is most likely to land}.\\

\begin{table}
\small
\begin{tabular}{@{}>{\raggedright}p{4.2cm} *{4}{>{\centering\arraybackslash}p{2.2cm}}@{}}
\toprule
\textbf{Directions} & \textbf{Routines (micro)} & \textbf{Infrastructures (meso)} & \textbf{Institutions (macro)} & \textbf{Ecosystems (inter-org)} \\
\midrule
Operational indicators & \Closed & \Closed & \OpenClosed & \Closed \\
Integrative frameworks & \Open & \Closed & \Closed & \Closed \\
Comparative \& cumulative & \Open & \Closed & \OpenClosed & \Open \\
Temporal/process markers & \Closed & \Open & \Closed & \OpenClosed \\
Normative/value-based & \Open & \OpenClosed & \Closed & \Closed \\
Governance and incentives & \Open & \OpenClosed & \Closed & \Closed\\
\bottomrule
\end{tabular}
\vspace{2mm}
\footnotesize Legend: $\Closed$ high opportunity; $\OpenClosed$ promising ; $\Open$ low opportunity.
\caption{Opportunity matrix locating research directions across organizational scopes (micro routines, meso infrastructures, macro institutions, inter-organizational ecosystems).}
\label{fig:resilience_interface_opportunity}
\end{table}

Reading the matrix along its rows brings these opportunities into focus. \emph{Operational indicators} (row 1) are most tractable at the micro, meso, and ecosystem scopes, where behavioural traces, infrastructure telemetry, and information-sharing signals are increasingly available for empirical work \autocite{bruneau2003framework,francis2014metric,panteli2017metrics,kuypers2018designing,gillard2023efficient}. The macro scope is currently rated as promising rather than high because authoritative cross-jurisdictional indicators remain harder to construct, although the harmonized incident classification introduced by DORA and the Basel Committee's principles for operational resilience offer a useful starting point for comparable macro-level measurement \autocite{bcbs2021operational,european2022regulation,nist2024csf}. \emph{Comparative and cumulative} research (row 3) shows a similar profile: micro-level case comparison and ecosystem-level studies are already well supported by existing methods \autocite{barasa2018what,haldane2021health,maillart2024computational}, while macro-level comparability across jurisdictions is improving but still partial \autocite{meerow2016defining,allen2018quantifying}. \emph{Temporal and process markers} (row 4) are particularly tractable at the infrastructural and ecosystem scopes, where event logs, recovery curves, and longitudinal incident records support process tracing \autocite{ives1995measuring,leslie2013response,aase2020resilience,krakovska2023resilience}. \emph{Governance and incentives} (row 6) are most directly addressable at the meso scope, where infrastructure governance, access design, and escalation pathways are concrete enough to be studied empirically \autocite{ostrom1990governing,lebel2006governance,boin2013resilient,finma2023operational}. Read together, the three artefacts---triangle, matrix, and opportunity table---thus specify what to study, where the mechanisms operate, and where empirical investment is most likely to be productive.\\

A related issue concerns cooperation under pressure. Time-critical crises rarely respect organizational boundaries. Response depends on whether actors can coordinate quickly across firms, infrastructures, and institutional settings. Yet this is exactly where trust asymmetries, activation delays, and incentive misalignments often undermine performance. Theories of collective action, modular governance, and platform-mediated collaboration therefore provide an important starting point, especially in digital threat environments where rapid mobilization is essential \autocite{ostrom1990governing,brauer2008compartmental,mahajan2022participatory}.\\

Empirical work on modular security platforms such as MISP is especially relevant here because it shows how structural design can enable or inhibit timely collaboration. Participatory simulations and controlled experiments add another layer by identifying thresholds for cooperative engagement and the points at which collaboration begins to fail structurally \autocite{skopik2016problem,gillard2023efficient}. Disaster-response research offers additional design principles for cross-organizational information sharing through IT systems, and many of those insights are transferable to adversarial digital environments \autocite{arnold2004informationsharing}. The depth of resilience as a substantive capability depends in no small part on how these institutional arrangements and cooperation dynamics are configured across scopes.

\subsection{Temporal Dynamics of Resilience Enactment}

A third research direction extends the enactment perspective by asking how resilience is formed and re-formed over time. The point is not simply to add a temporal layer to the studies discussed above, but to examine how enacted capacities are built, stabilized, weakened, or redirected through experience, feedback, and institutional choice. On this view, resilience is neither a fixed organizational attribute nor something secured by resource availability alone. It is a contingent and relational capacity whose durability depends on whether earlier adaptations are reinforced, remembered, and carried forward, or instead neglected and allowed to erode \autocite{rutter1987psychosocial,luthar2000construct,ungar2021modeling}. This temporal dimension is comparatively underdeveloped in the IS literature, which still tends to focus on discrete events or isolated responses. Organizations do not confront each disruption from a blank slate: they carry residues of earlier crises---routines, assumptions, temporary fixes, and interpretive habits that, in \citeauthor{lanzara1999transient}'s analysis of post-disaster organizing, function as transient yet consequential scaffolds for later action \autocite{lanzara1999transient}. \textcite{westrum2004typology}'s typology of generative, bureaucratic, and pathological information cultures provides a complementary anchor for thinking about whether such residues become durable organizational learning or remain inert. Any persuasive account of resilience therefore has to explain not only how organizations respond in the moment, but also how adaptive capacity is retained, revised, or lost across successive episodes of disruption.\\
\\
Useful foundations for this line of inquiry already exist in systems theory and ecology. \citeauthor{holling1973resilience}'s work on adaptive cycles draws attention to recurring phases of growth, conservation, release, and reorganization, making clear that resilience is tied to sequences and transitions rather than to any single moment of adjustment \autocite{holling1973resilience,gunderson2002panarchy,folke2010resilience}. \citeauthor{wiener1948cybernetics}'s cybernetics places feedback at the center and shows how systems may be stabilized through corrective feedback or pushed toward instability through reinforcing feedback \autocite{wiener1948cybernetics}. More recent work on basin stability and return-time dynamics formalizes the same intuitions for nonlinear systems \autocite{gao2016universal,krakovska2023resilience}. Taken together, these traditions suggest that resilience depends not only on an organization's immediate ability to cope with disruption, but also on whether effective responses are reinforced, whether lessons become institutionalized, and whether failure produces revision rather than repetition \autocite{folke2004regime,allen2018quantifying}. Under some conditions resilience accumulates over time; under others it decays through forgetting, inertia, or maladaptation; and under the antifragile conditions described by \textcite{taleb2012antifragile}, controlled exposure to stressors may even compound adaptive capacity, provided the supporting learning routines remain durable.\\

Studying these dynamics calls for methodological pluralism. Longitudinal case studies are needed to follow how organizations adapt across repeated disruptions rather than single events. Process tracing can then identify the pathways through which learning practices, governance arrangements, and changes in technical systems shape later outcomes. Computational models offer complementary leverage by making it possible to examine how different rates of reinforcement, forgetting, or knowledge transfer generate different long-run resilience trajectories. Field evidence from Appendix~\ref{appendix:C} (theme iv) underscores why this temporal lens matters in practice: respondents repeatedly described post-incident reviews as conducted but rarely institutionalized (\textit{``The same issues resurface again and again because lessons learned are not retained''}), making the half-life of organizational learning a tractable and high-value object of empirical study. Taken together, these approaches would allow IS research to move beyond largely static conceptions of resilience and instead analyze it as a socio-technical capability that can strengthen, deteriorate, or mutate as organizations are repeatedly disturbed.

\subsection{Toward an Integrative Enactment-Scope Research Agenda}

Taken together, these research directions point toward a broader agenda: resilience can be studied most fruitfully by linking mechanisms of enactment to the organizational scopes within which they operate. Building on the contributions of fragmented perspectives, the field is now positioned to develop an integrated socio-technical account that relates behavioral processes, institutional conditions, and temporal development without collapsing them into a single level of analysis. This ambition is consistent with systems and ecological traditions that treat resilience as an emergent effect of feedback, interdependence, and adaptation \autocite{wiener1948cybernetics,holling1973resilience}. For IS research, that means building explanations that connect micro-level action, meso-level organizational and infrastructural dynamics, and macro-level governance constraints. It also means developing indicators that complement compliance reporting with finer-grained traces of adaptation as it unfolds: changes in response variability, shifts in coordination under stress, and the reconfiguration of technical and organizational arrangements during disruption \autocite{aerts2018integrating}. Anchoring resilience in observable behavior is the bridge that helps the conceptual ambition of the field meet analytical precision.\\
\\
At the same time, resilience does not emerge in a vacuum. The institutional and organizational environments in which actors operate shape what forms of adaptation are feasible in the first place. Rules, incentives, authority structures, and escalation pathways do not simply constrain action after disruption; they structure the range of possible responses before disruption occurs \autocite{ostrom1990governing,ungar2021modeling}. Existing field-based and computational studies already show that these conditions matter, often appearing as friction, delay, coordination breakdown, or reduced strategic flexibility. The harder task now is to turn these scattered insights into analytical frameworks that can travel across settings without losing explanatory force.\\

Figure \ref{fig:resilience_interface_opportunity} makes this research space visible by locating opportunities across micro behaviors, meso infrastructures, macro institutions, and inter-organizational ecosystems. Read together, Figures \ref{fig:triangle} and \ref{fig:resilience_interface_opportunity} define the central logic of the Enactment--Scope framework. The triangle specifies how resilience moves from articulated organizational commitment to enacted practice. The opportunity matrix shows where those mechanisms can be studied and where intervention is most likely to matter. Advancing this agenda will require research programs that combine theoretical ambition with methodological range, including computational modeling, longitudinal case analysis, and field-based validation. Used together, these approaches offer a more credible path toward cumulative and transferable knowledge about organizational resilience \autocite{teece2007explicating,gillard2023efficient,kalisch2024neurobiology}. More importantly, they provide a stronger basis for IS theorizing about digital resilience by showing how adaptive capacity emerges through the interaction of technological architectures, organizational practices, and institutional environments. On this view, resilience is best understood as more than a strategic statement or a checklist: it is an enacted, distributed, and revisable socio-technical capability for maintaining and reorganizing coordinated action under persistent uncertainty \autocite{weick2011managing,woods2015four}.
<<<<<<< HEAD
\section{Implications for Research \& Practice}
=======
\section{\MakeUppercase {Implications for Research \& Practice}}
>>>>>>> 5962c492beeeeff7385c8b177888710a79b39489
\label{sec:ImplicationsResearch}
\renewcommand{\baselinestretch}{2.0}\normalsize

<<<<<<< HEAD
The three gaps identified in this review -- conceptual ambiguity, limited empirical grounding, and methodological fragmentation -- do not point to isolated weaknesses. Together, they highlight a wider issue with the way resilience is studied in digitally mediated environments. The issue is not just that the literature is incomplete. Rather, resilience is still too often invoked in general terms, measured inadequately, and only operationalised after disruption has already occurred. The following implications take a different view. They treat resilience as a socio-technical capability that must be defined more precisely, observed in action and linked to effective interventions.\\
\\
Firstly, resilience should not be considered a fixed organisational attribute. It is not merely a residual label attached to organisations that appear to cope well after disruption, either. It is better understood as an emergent socio-technical capability. This is important because the literature often moves too easily between resilience, robustness, recovery and agility\autocite{woods2015four,duchek2020organizational}. The overlap is real, but so are the differences. In digital settings especially, resilience is not confined to one asset, one team, or one managerial disposition. Rather, it emerges through interaction between people, technical systems, routines, escalation pathways and governance arrangements. A good example is cyber incident response, where it isn't about one particular resource but rather how well the analysts, security software, and other reporting protocols, along with other channels for information sharing, can collaborate successfully within time constraints. Resilience framed in such a manner will give IS researchers a better starting point for their analyses, while facilitating a more productive interaction between IS research and dynamic capability studies.\autocite{henfridsson2013generative,boh2023building}, while remaining compatible with systems and ecological traditions \autocite{holling1973resilience,wiener1948cybernetics}.\\
\\
A second implication follows from this almost immediately. Resilience must be made observable. Otherwise, it remains just a slogan. Behavioural and computational indicators are important not because they make resilience appear more scientific, but because they make the concept empirically accountable. Plausible candidates include coordination speed, stakeholder support, community ties, the capacity to reconfigure infrastructures, and variation in response under stress \autocite{mayer2021putting,hollnagel2013resilience,gillard2023efficient}. Once specified, such indicators can be studied through agent-based simulation, network analysis, or probabilistic modelling \autocite{helbing2012agent}. This is especially important in digital spaces, where adaptation is spread out, needs to happen quickly, and is often only partially visible. A simulation model can show how changes in how quickly information is shared or how weak incentives for collaboration can change how people work together during cyber incidents. That is the point. When you can see something, you can compare resilience. It can also be questioned and changed.\\

This also means we spend less time treating compliance artifacts as evidence of resilience. Although audits, certifications, and maturity assessments are not useless, they are weak indicators of adaptive capacity in practice.  \autocite{skopik2016problem}. They inform us of the organizations that claim to be prepared to do so. They divulge scant information regarding the ramifications of infrastructure failure, the dissolution of coordination, or the transformation of circumstances midway. Field-based work is therefore essential. Longitudinal case studies, trace data, and ethnographic approaches are important here. Research in healthcare and public health already shows that resilience can be evident in everyday practice while remaining only partially institutionalized \autocite{kruk2017building,barasa2018what}. The same opportunity exists in digital infrastructures, where incidents and near-misses reveal how organisations improvise, adapt, stall or fail over time. In some cases, disruption may trigger processes of renewal reminiscent of Schumpeter's 'creative destruction', in which routines or infrastructures that no longer fit prevailing conditions are replaced by more adaptive ones  \autocite{schumpeter1942capitalism}. This should not be romanticised. However, it does reinforce the broader point that resilience becomes analytically meaningful when studied in action rather than inferred from documentation.\\
\\
Another implication is that adaptive capacity does not belong to only one level of analysis. Rather, it unfolds across micro-level routines, meso-level infrastructures, macro-level institutions, and inter-organizational ecosystems. This is one reason why single-level explanations often feel inadequate. A cyberattack may begin with analyst decisions at the micro level, spread through infrastructural dependencies at the meso level, and then be redirected or constrained by regulation, liability, or cross-organizational coordination at broader scales. The mechanisms are distributed. IS research is unusually well placed to deal with that because it can combine process tracing, computational modeling, and comparative case analysis without forcing the problem into a single analytical register \autocite{ostrom1990governing,vogus2007organizational,folke2010resilience}. That is a real advantage of the field. It also creates room for indicators -- such as coordination speed, recovery sequence, or adaptive reconfiguration -- that travel across contexts without collapsing into empty abstraction.\\
\\
There is also a practical implication that should be stated plainly. Better concepts and indicators won't matter much if they don't influence how organizations prepare, make decisions, and coordinate. To that end, the field of resilience research must evolve to become more managerially useful without succumbing to managerial clichés. This requires design principles, benchmarks, early warning systems, and governance mechanisms that serve a purpose beyond mere decoration of annual reports \autocite{weick2011managing,lengnick2011developing,calder2021iso}. The risk is obvious: Organizations can demonstrate resilience long before they operationalize it. Metrics theater is not rare. Dashboards exist, escalation paths are documented, and the language of resilience is prevalent. Yet, adaptive capacity remains thin when pressure actually arrives. Therefore, linking measurement to governance is not optional. In practice, this may involve real-time monitoring of the adaptive response, escalation protocols that clarify authority during disruption, or governance arrangements that assign coordination responsibilities before a crisis begins.\\
\\
When considered as a whole, these findings suggest a need for a more rigorous research agenda in the field of IS. The field does not need more generic endorsements of resilience or another round of compliance-friendly abstraction. To achieve this, it is essential to develop more sophisticated concepts, strengthen the empirical foundation, and devise methods that can elucidate the process by which adaptive capacity emerges in the face of ongoing disruption. For scholars, that creates room for more cumulative theory. It also creates room for better evidence. The conversation is shifted for practitioners from symbolic commitment to capabilities that can be observed, tested, and revised. Reframed in these terms, resilience becomes less a rhetorical aspiration than a grounded socio-technical capacity for operating under strain in digitally mediated environments.
=======
Regulated sectors---finance above all---already run on a mature digital industry backbone. Business continuity plans (BCPs) and business continuity management systems aligned with ISO~22301 \autocite{calder2021iso}; business impact analyses; recovery-time and impact-tolerance objectives for important business services under BCBS operational-resilience principles and national regimes such as FINMA guidance and DORA \autocite{bcbs2021operational,european2022regulation,finma2023operational}; software supply-chain dependency mapping \autocite{shen2024understanding}; ICT continuity and response-and-recovery policies; annual backup-and-restore testing; crisis-communication and third-party dependency mapping; and threat-led penetration-testing programs (TIBER-EU, CBEST) \autocite{ecb2018tiber,boe2021cbest}---together constitute state-of-the-art preparedness, not a missing control baseline. On that backbone, organizations add SOCs, Security Information and Event Management (SIEM) telemetry, and sector exercises---then still struggle to show boards how any of it evidences adaptive performance when disruption departs from script (\autocite{rota2024}; Appendix~\ref{appendix:C}). BCPs are invoked; cross-unit handoffs stall; ticket histories and after-action findings rarely alter governance rhythms. Resilience is declared in every committee; enacted adaptation remains hard to demonstrate. The deeper problem is conceptual: resilience is still treated as a validated achievement---maturity attainment, plan certification, audit clearance---when it is more accurately a continuous journey of learning and exploration within and across organizational boundaries. Research and practice must shift from certifying possession of controls to validating that journey itself---whether sensing, reconfiguration, and coordination actually improve as threat and disruption conditions evolve.\\
\\
Section~\ref{sec:RD} turns that reorientation into three research imperatives, scoped in Figure~\ref{fig:enactment_scope} and prioritized in Table~\ref{fig:resilience_interface_opportunity}. \textbf{Define} what counts as adaptive performance---at infrastructure cascade, internal coordination, and cross-firm activation---before BCP templates and audit scores substitute for it \autocite{hillmann2021organizational}. \textbf{Measure} it from digital traces organizations already generate: SOC and SIEM logs, BCP test and invocation records, recovery timelines, coordination metadata, ISAC records. Start with explainable quantitative methods; add machine-learning and explainable-AI layers where volume demands it---orienting commercial analytics platforms toward governance-facing indicators rather than detection tickets that never leave the security function \autocite{francis2014metric,gillard2023efficient}. \textbf{Govern} it: bind those indicators to BCP escalation paths, vendor recovery, board review, and cooperative activation across partners---and to the transient governance that disruption triggers, instrumenting how decision authority migrates and hands back rather than only the steady-state controls---not slide decks assembled after incidents close \autocite{boin2013resilient,skopik2016problem,ecb2018tiber}. Where single-firm logs cannot reveal path-dependent cascade across cloud and supplier dependencies, simulation must supply the counterfactual ensemble \autocite{dawson2011agentbased,kwakkel2016comparing}. IS research sits at the hinge: the traces, infrastructures, and cross-boundary workflows through which resilience is actually produced are already digital.

The implications below translate that reorientation for IS scholarship and for the organizations whose infrastructures, routines, and partnerships now constitute the primary site of adaptive performance under disruption.\\
\\
Firstly, the field must confront an oddly asymmetric blind spot. Management has long theorized and instrumented adaptation to market conditions---gradual repositioning under stable demand and sharp strategic pivot under shock \autocite{teece2007explicating,duchek2019organizational}---yet still frames security and operational resilience chiefly as extensions of risk management and compliance. Continuously adapting, forecasting, and measuring security postures and coordination capabilities in response to evolving threat landscapes remains comparatively under-theorized and under-instrumented. Because resilience is annexed to security and risk management in both research and practice, it largely misses the perspective shift its own logic demands: operating under disruptive conditions is less a control problem than an ongoing business-adaptation problem. Resilience is therefore best understood not as a fixed organisational attribute, nor only as a label assigned to organisations that appear to cope well after disruption, but as an emergent socio-technical capability enacted across those scopes and mechanisms---a journey whose progress, not its declaration, requires validation. The literature still moves between resilience, robustness, recovery, and agility without distinguishing them sharply \autocite{woods2015four}. In digital settings especially, resilience is not confined to one asset, one team, or one managerial disposition; it emerges through interaction between people, technical systems, routines, escalation pathways, and governance arrangements. Cyber incident response illustrates the point: the outcome depends less on any single resource than on how well analysts, security software, reporting protocols, and information-sharing channels collaborate within time constraints. Framed in this way, resilience provides IS researchers a clearer starting point for analysis and a more productive interaction with dynamic capability studies \autocite{henfridsson2013generative,boh2023building} and digital-transformation research \autocite{vial2019understanding,hanelt2021systematic,wessel2021unpacking}, while remaining compatible with systems and ecological traditions \autocite{wiener1948cybernetics,holling1973resilience}.\\
\\
A second implication follows directly. Validating the resilience journey requires making adaptive trajectories observable---not only snapshots of preparedness at a point in time. Behavioural and computational indicators serve that purpose: not by making resilience appear more scientific, but by giving practitioners and researchers a shared evidentiary basis for whether learning and reconfiguration are improving over successive disruptions, exercises, and near-misses. Plausible candidates include coordination speed, stakeholder support, community ties, the capacity to reconfigure infrastructures, and variation in response under stress \autocite{paries2017resilience,gillard2023efficient}. Once specified, such indicators can be studied through agent-based simulation \autocite{helbing2012agent}, network analysis \autocite{panteli2017metrics}, or probabilistic modelling \autocite{kruk2017building}. This is especially valuable in digital spaces, where adaptation is distributed, time-critical, and only partially visible. A simulation model can show, for instance, how variations in the speed of information sharing or the strength of collaboration incentives shape how people coordinate during cyber incidents. Once made visible, the journey can be compared across cases, examined more carefully, and improved deliberately.\\
\\
A further implication cuts across all three imperatives: artificial intelligence is becoming the connective layer through which they can be pursued at scale, in research and in practice. On the measurement and decision side, explainable AI can fuse the heterogeneous signals resilience leaves behind---SOC and SIEM telemetry, coordination metadata, exercise after-action data, and cross-firm sharing records---into decision-ready indicators, provided its reasoning remains legible to operators, boards, and regulators rather than collapsing into opaque risk scores \autocite{berente2021managing,gillard2023efficient}. On the anticipatory side, reinforcement-learning and related generative methods can search large spaces of adversarial and operational contingencies to surface \emph{robust} future scenarios---the regimes a strategy must be able to withstand---whose mastery is precisely what adaptation under deep uncertainty requires \autocite{kwakkel2016comparing,lempert2019robust}. Framed this way, AI does not displace the construct, indicators, and governance developed above; it integrates them, turning dispersed traces into foresight and foresight into rehearsed response. The same opacity that makes such models powerful, however, makes explainability a precondition rather than an enhancement if their outputs are to carry weight in governance.\\

This also implies extending---rather than replacing---the role of compliance artifacts and continuity planning in resilience assessment. Audits, certifications, business continuity plans, maturity assessments, and BCP test records provide essential accountability and have advanced organizational governance considerably; their range can now usefully be complemented by direct measures of adaptive performance \autocite{skopik2016problem,calder2021iso,bcbs2021operational}. They tell us a great deal about how organizations articulate preparedness; combining them with field-based observation captures more of what happens when infrastructure fails, coordination is reorganized, or circumstances shift midway. Longitudinal case studies, trace data, and ethnographic approaches are particularly well suited to that complement. Research in healthcare and public health already shows that resilience is often evident in everyday practice well before it is fully institutionalized \autocite{kruk2017building,barasa2018what}. A similar opportunity exists in digital infrastructures, where incidents and near-misses offer rich material on how organisations improvise, adapt, recover, or fall short over time. In some cases, disruption triggers processes of renewal reminiscent of Schumpeter's 'creative destruction', in which routines or infrastructures that no longer fit prevailing conditions are replaced by more adaptive ones. Without romanticising disruption, this reinforces the broader point that resilience becomes analytically richest when studied in action alongside the documentation that organizations also produce.\\
\\
Another implication is that adaptive capacity does not belong to only one level of analysis. It unfolds across micro-level routines, meso-level infrastructures, macro-level institutions, and inter-organizational ecosystems. This is one reason why single-level explanations capture only part of the picture. A cyberattack may begin with analyst decisions at the micro level, spread through infrastructural dependencies at the meso level, and then be redirected or constrained by regulation, liability, or cross-organizational coordination at broader scales. The mechanisms are distributed. \textcite{bhamra2011resilience}'s synthesis of the resilience literature already pointed to this multi-level architecture, and subsequent work on the transient governance that disruption induces (Section~\ref{sec:RD}), information cultures that convert near-misses into learning \autocite{westrum2004typology}, and polycentric governance of common resources \autocite{ostrom1990governing} provides complementary entry points. IS research is unusually well placed to engage with that complexity because it can combine process tracing, computational modeling, and comparative case analysis without forcing the problem into a single analytical register \autocite{vogus2007organizational,folke2010resilience}. That is a real advantage of the field. It also creates room for indicators that remain anchored in concrete organizational practice as they travel across contexts.\\
\\
There is also a practical implication worth stating plainly. Better concepts and indicators have most value when they shape how organizations prepare, make decisions, and coordinate. The field of resilience research can therefore continue to grow into a more useful body of work, building on the substantial managerial and regulatory effort already devoted to design principles, benchmarks, early warning systems, and governance mechanisms \autocite{calder2021iso,bcbs2021operational,nist2024csf}. A practical caution accompanies this progress: articulated resilience can outpace its operationalization, and many organizations are already aware of the gap. Field interviews \autocite{rota2024} capture this tension directly: boards request dashboards and KPIs, yet defining and resourcing indicators that boards will actually use often remains unresolved (Appendix~\ref{appendix:C}). Dashboards exist, escalation paths are documented, and the language of resilience is widely shared; the next step is to ensure that these structures connect tightly to adaptive capacity when pressure arrives. Linking measurement to governance is therefore central. In practice, this may involve real-time monitoring of the adaptive response \autocite{kuypers2018designing,gillard2023efficient}, escalation protocols that clarify authority during disruption \autocite{westrum2004typology,boin2013resilient}, or governance arrangements that assign coordination responsibilities before a crisis begins \autocite{ostrom1990governing,arnold2004informationsharing}. Some of these arrangements, in the spirit of \textcite{taleb2012antifragile}, can also be designed to generate productive learning from controlled exposure to stressors, as the threat-led testing regimes formalized by CBEST and TIBER-EU illustrate \autocite{ecb2018tiber}.\\
\\
Considered as a whole, these findings invite a more focused resilience research agenda in the field of IS. Building on the conceptual ground already laid by management, organization, and IS scholars, the next step is to refine the constructs, strengthen the empirical foundation, and develop methods that trace how adaptive capacity evolves---not merely whether controls were present---across successive disruptions. For scholars, this creates room for more cumulative theory and richer evidence on the resilience journey. For practitioners, the conversation can shift from claiming preparedness to demonstrating that adaptation is learning, measurable, and governable alongside the strategic commitments organizations have already made. In these terms, resilience moves from an organizing principle to a grounded socio-technical capacity for operating under increased strain in digitally mediated environments \autocite{majchrzak2016designing}.
>>>>>>> 5962c492beeeeff7385c8b177888710a79b39489

\section{\MakeUppercase {Conclusion}}
\label{sec:Conclusion}
\renewcommand{\baselinestretch}{2.0}\normalsize


Resilience is a serious organizing concept in management and information systems research, increasingly central to how organizations articulate their response to disruption. Regulated sectors have built a mature preparedness backbone---continuity planning, operational-resilience regimes, threat-led testing---and the field has produced substantial conceptual, regulatory, and methodological work. Yet researchers and practitioners still struggle to show that resilience is more than declared: adaptive performance under disruption, and improvement across successive events, often remains harder to evidence than compliance itself \autocite{rota2024,skopik2016problem,woods2015four}.\\
\\
This review identifies three interlocking deficits---conceptual plurality, limited empirical validation, and methodological fragmentation---and responds with the Enactment--Scope framework and the research agenda developed in Sections~\ref{sec:Enactment-Scope}--\ref{sec:RD}. The implications for practice are sharper than accumulating further controls. Resilience is better treated as a continuous journey of learning and exploration within and across organizational boundaries than as a validated achievement certifiable through maturity scores or audit clearance. What requires validation is that journey itself: whether sensing, reconfiguration, and coordination improve as threat and disruption conditions evolve. Management already knows how to instrument adaptation to market conditions \autocite{teece2007explicating,duchek2019organizational}; the analogous task for security postures and coordination capabilities under evolving threats remains comparatively under-developed.\\
\\
For IS research, three imperatives follow: \textbf{define} adaptive performance before templates substitute for it; \textbf{measure} it from digital traces organizations already generate; and \textbf{govern} it through escalation, review, and cooperative activation rather than post-incident reporting alone. By advancing construct clarity, trace-based indicators, simulation under deep uncertainty, and governance designs that connect measurement to decision rights, information systems research is well positioned to help organizations validate adaptive capacity in digitally mediated environments---moving resilience from organizing principle to enacted, measurable, and revisable socio-technical capability.
% References
%\newpage
\printbibliography

\input{sections/9_appendices}

\end{document}